\section{RMSNorm 参数补充}

\subsection{RMSNorm 的结构}

RMSNorm（Root Mean Square Normalization）是 LLaMA 2 引入的归一化方法，比 LayerNorm 更简洁——只做缩放，不做均值中心化：

$$
\text{RMSNorm}(x) = \frac{x}{\sqrt{\frac{1}{d}\sum_{i=1}^d x_i^2 + \epsilon}} \odot \gamma
$$

其中 $\gamma \in \mathbb{R}^{d_{\text{model}}}$ 是可学习的缩放参数。

\subsection{RMSNorm 的分布}

每个 Transformer Layer 包含 \textbf{2 个} RMSNorm：
\begin{enumerate}
  \item GQA Attention 前的 Norm
  \item FFN/MoE 前的 Norm
\end{enumerate}

加上模型最后的 \textbf{1 个} Final Norm，总计：

$$
P_{\text{RMSNorm}} = (2 \times L + 1) \times d_{\text{model}}
$$

对于 Dense 32B：$(2 \times 64 + 1) \times 5120 = 660{,}480$ 个参数。虽然相对于 32B 可忽略不计，但概念上不应遗漏。

\subsection{本章小结}

RMSNorm 是每层的「标配」，参数量很小但不可或缺。它保证了每层输入的数值稳定性，是训练深层 Transformer 的关键组件之一。

